\section{Related work}
\label{sec:related}

This section covers the work the rest of the report builds on, and one class of
work it deliberately declines. The background the main argument does not depend
on is collected in \S\ref{app:related}, the evaluation-methodology literature is
taken up in \S\ref{sec:protocol}, and \S\ref{app:rel-declined} states, for each
declined method, why it is declined. Absence is a design decision in this study,
not an omission, and it is recorded as one.

Three of this report's results sit against specific literatures and are placed
here rather than left implicit. The optimiser's-curse result of
\S\ref{sec:search-curse} belongs to the determinization line, where an unguarded
argmax over sampled worlds is a known-in-principle failure that is rarely
quantified. The belief result of \S\ref{sec:belief-exactworse} --- a
policy-aware approximation beating the exact posterior --- belongs to the
trick-taking inference line, which is the only literature that has measured the
same trade-off. And the shared common-knowledge object of
\S\ref{sec:mech-knowledge} is the common-information reduction, restricted to
the one case in which it is cheap.

\subsection{Fish and Literature}
\label{sec:related-fish}

The published material on this game is human and informal, and neither of its
two substantial sources is quantitative or describes a program. McLeod's Pagat
entry \cite{pagat} is the canonical rules source and records several points on
which dialects differ; ours is specified in \S\ref{sec:rules} and listed in full
in \S\ref{app:dialect}. Develin's chapter \cite{develin} is the most detailed
strategy writing we located, describes the same dialect, and supplies two
observations this design uses: an ask teaches the opponents as much as it
teaches a teammate, and a player may not declare a half-suit they hold entirely.
Secondary strategy pages \cite{dorsa,strategy-site,wikiliterature} restate
blackballing, asking the asker, and an endgame delegation heuristic; they are
evidence about how the game is played and talked about, not sources of
quantitative claims. \S\ref{app:rel-informal} records where the two substantial
sources disagree.

Computational prior art remains close to absent. The v0.4 study's searches found
no academic work on this game, no public repository containing search, planning
or a working learner for it, and no coverage in RLCard \cite{rlcard}; the v0.6
literature refresh re-ran that question against 2025--2026 work and found
nothing that changes it. The one agent located, Somani's \texttt{literature}
\cite{somani,somani-server}, plays the four-player, 48-card variant, which does
not raise the six-player allocation problem, and publishes no strength numbers;
a closed-source Android application advertises bots with no stated methodology
\cite{playstore}. Every novelty statement in this report therefore takes the
form ``we found no prior published work that \dots'' and should be read against
that scope. One 2026 item is worth naming because it may become the external
calibration this line of work lacks: Valet packages twenty-one traditional
imperfect-information card games in a common description language with measured
branching factors and durations \cite{goadrich2026valet}, though whether any of
the twenty-one is a team game is not stated on its abstract page and we did not
fetch the full text, so nothing here is compared against it. The three prior
FishBot reports \cite{fishbot03,fishbot04,fishbot05} are the record this study
extends and, in \S\ref{sec:corrections}, corrects.

\subsection{Determinization, and why double-dummy analysis fails here}
\label{sec:related-determinization}

Determinization, or perfect-information Monte Carlo --- sample hidden cards
consistently with the public history, solve the resulting perfect-information
game, and play the move with the best average value --- is the standard recipe
for hidden-hand card games \cite{levy,ginsberg-gib}, and is unsound however many
worlds are drawn. Frank and Basin named the two mechanisms: \emph{strategy
fusion}, in which the searcher plays a different strategy in each sampled world
although a real player must commit to one strategy per information set, and
\emph{non-locality}, in which a node's value depends on parts of the tree
outside its own subtree \cite{frank-basin-ai,frank-basin-aaai}. Long, Sturtevant,
Buro and Furtak characterise when the method nevertheless performs well, through
leaf correlation, bias and a disambiguation factor \cite{long-pimc}; Information
Set Monte Carlo Tree Search replaces independent trees with a single tree over
information sets \cite{cowling-ismcts,goodman-ris}; the $\alpha\mu$ family backs
up Pareto fronts of world-vectors instead of scalar averages, repairing both
mechanisms, and reports Bridge 3NT declarer play at 60.2\% to 63.0\% with twenty
worlds \cite{cazenave2019alphamu}; and knowledge-based paranoia search reasons
about the worst case over deduced-consistent worlds rather than the average,
reporting over 1,000 points on the extended Seeger scale in Skat
\cite{edelkamp2021paranoia}. Its runtime is not stated on the abstract page, so
no compute cost is attributed to it anywhere in this report.

Fish breaks the method for a reason none of that literature addresses. Its
perfect-information game is not merely fused, it is degenerate: a clairvoyant
player names only cards the target actually holds, so never fails an ask, never
loses the turn, and declares each half-suit correctly as soon as their team
holds all of it. Almost every reachable position carries the same
perfect-information value, and averaging that value over deals sampled from the
public history collapses to choosing the ask most likely to hit --- a one-ply
heuristic blind to the information an ask leaks, the information a refusal
supplies, and the option value of postponing a declaration.
$\alpha\mu$ repairs unsoundness, not degeneracy, so it does not rescue the
construction here. \S\ref{app:rel-determinization} gives the collapse in full.

What \S\ref{sec:search} adds is that the binding problem in this game is not the
abstract failure but the statistic. One determinization's value is the final
half-suit differential, with standard deviation \vsixLeafSd{} --- an order of
magnitude larger than the differences between the leading candidates --- so an
argmax over $D$ sampled means selects on residual noise, and does so more
confidently the more candidates it is offered. That is the optimiser's curse
rather than strategy fusion, and it is measured here at \vsixCurseGap{} points
by a control that forces the same code to return the blueprint's own pick
(\S\ref{sec:search-curse}). This report therefore uses determinization only as a
research probe, behind the paired deviation guard of \S\ref{sec:search-falsifier},
and reports the result including the configurations that land at zero
(\S\ref{sec:search-deadask}). The transferable claim is the guard, not the
search.

\subsection{Search in imperfect-information and cooperative games}
\label{sec:related-search}

States in an imperfect-information game do not have well-defined values, so
depth-limited search needs something at the leaf that represents the opponent's
ability to choose. Brown, Sandholm and Amos supply it: let the opponent select
among $k$ continuation strategies at the depth limit, each yielding its own
leaf-value vector, and require the searcher to be good against that choice ---
demonstrated at master level in heads-up no-limit hold'em on a four-core CPU and
16\,GB \cite{brown2018depthlimited}. That construction is the one piece of this
literature this study uses as method rather than context: it is the multi-valued
leaf of \S\ref{sec:mech-leaves} and the prerequisite for
\S\ref{sec:search}. The soundness argument is stated for two-player zero-sum
games, and Fish's team-level game presents three uncorrelated opponent seats at
the leaf, so nothing is inherited. It is adopted here as a robustification
heuristic with a measurable consequence, and audited rather than assumed
(\S\ref{sec:mech-exploit}).

Two lines address the other half of the problem, which is what a search may
prune. Zhang and Sandholm solve subgames over low-order knowledge rather than
the full common-knowledge closure, which is the published answer to a closure
too large to enumerate --- and they prove that the naive form can \emph{increase}
exploitability \cite{zhang2021subgamesolvingwithoutck}. That warning applies
directly to the live-ask gate this policy already ships, which is a
knowledge-limited pruning rule that has never been audited, and it is why
\S\ref{sec:mech-exploit} was built before any search work rather than after it.
Subgame solving in adversarial team games builds a gadget over the team belief
DAG and runs in polynomial time given a best-response oracle \emph{when the
blueprint is sparse} \cite{zhang2022teamsubgame}; a Fish blueprint is not
sparse, so that escape hatch does not open here, and \S\ref{app:rel-declined}
records this as a verified non-transfer rather than an unexamined gap.

In cooperative games the constraint is different and it transfers verbatim.
SPARTA improves a blueprint purely at test time, 24.08 to 24.61 in Hanabi, and
distinguishes single-agent search --- which assumes every other agent strictly
follows the blueprint --- from multi-agent search, in which all agents run the
\emph{same} common-knowledge procedure \cite{lerer2020sparta}. The moment a
second teammate also searches, the first teammate's belief over its partner is
computed under a policy that partner is no longer playing, and the improvement
guarantee is void. This is the Dec-POMDP common-knowledge caveat, and it is the
constraint most easily violated in silence. Fish satisfies it cheaply: all card
movement is public, so any \emph{deterministic} procedure that reads only public
state and its own hand is automatically common knowledge among the three
teammates, with no shared seed and no communication. The qualifier is
load-bearing --- a randomised tie-break loses the property unless its
randomisation is seeded from the public history, which is a one-integer decision
with a real correctness consequence (\S\ref{sec:mech-tiebreak}).

The remainder of this literature is cited as context. The full-solver line
\cite{moravcik2017deepstack,brown2020rebel,schmid2023sog} established that
test-time re-solving over public belief states is principled rather than a hack,
and ReBeL's safety result is the licence under which any test-time search here
operates; the algorithms are declined, because each requires iterating or
learning over a closure whose cardinality at deal time is
$45!/(9!)^5 \approx 1.9 \times 10^{28}$ deals from one seat. Update equivalence
recasts search as replicating a last-iterate-convergent update locally, needing
no public-information enumeration at all, and reports matching or exceeding
public-belief-state search in Hanabi at two orders of magnitude less search time
\cite{sokota2024updateequivalence}; its improvement guarantee is for
common-payoff games, which Fish is not, so it informs the shape of the mechanism
set without licensing it. Learned Belief Search buys 55--91\% of exact search's
benefit at 4.6--42$\times$ less compute \cite{hu2021lbs} and is simply redundant
here, because this game's exact belief is already computed and oracle-validated
(\S\ref{sec:belief-count}). Online planning by policy fine-tuning
\cite{fickinger2021rlfinetuning} requires gradient updates at decision time and
does not transfer to a CPU-only engine with no network. The depth-limited
counter-strategy line \cite{milec2021cdr,milec2025abd} carries one warning this
study honours --- robust-response methods do not compose with limited-look-ahead
solving off the shelf --- and the learned-model and gadget-refinement work
\cite{kubicek2025lamir,kubicek2026refinements} is recorded for its meta-lesson
rather than its machinery: when a search subproblem has many near-equivalent
optima, which one is selected is not a tie-break, and refinement is reported at
more than 50\% of exploitability. That is the same observation
\S\ref{sec:diag-ties} makes about an argmax resolved by array order.

Finally, every method above is published with a caveat that it can increase
exploitability, so none of them can be evaluated by win rate alone. Local best
response gives a cheap lower bound: a small greedy search, at most two actions
deep, showed every entrant in a 2016 poker competition to be more than
3,180\,mBB/h exploitable \cite{lisy2017lbr}. Its two documented pitfalls are
carried into \S\ref{sec:results-lbr} --- restricting where the response may
begin deviating understates exploitability, and a negative score proves nothing
--- and the figure this report publishes is a lower bound within its response
class, not an exploitability number.

\subsection{Policy-aware inference}
\label{sec:related-inference}

Skat engines enumerate the deals consistent with the public record and reweight
them by a model of what each player would have done --- a supervised per-card
location model \cite{solinas-supervised}, or reach probabilities taken from a
policy \cite{rebstock-inference} --- and a dedicated inference module was
reported worth $+217$ tournament points per 36 games to the engine Kermit on its
own \cite{buro-kermit}. The policy prior of \S\ref{sec:belief-prior} is the same
idea in a much simpler form: a fitted soft reweighting of the rules posterior,
with two parameters and no learned opponent model.

This literature is where the central result of \S\ref{sec:belief} belongs, and
it is the only literature that has measured the same trade-off. What is new here
is the margin and its direction. In this game the exact conditional posterior is
computable in closed form and cheap (\S\ref{sec:belief-count}), so the usual
excuse for approximating --- that exactness is unaffordable --- does not apply;
and yet substituting the exact posterior into an otherwise identical policy
\emph{loses}, at two independent seed banks. The exact path resolves more
distinctions and predicts hits worse, and the reason is structural rather than
numerical: the block construction consumes the deduction state alone and has no
parameter through which behaviour could enter, while the approximate path is
policy-aware by construction. Exactness and behaviour-awareness are not
available in the same object here, and \S\ref{sec:mech-tiebreak} is built on
that split. The report is careful about what this does and does not show: the
weights were fitted for the approximate belief, so it is a substitution result
rather than a verdict on exact inference (\S\ref{sec:limitations}).

Two adjacent lines are named and not used. Belief modelling in cooperative games
--- the Bayesian Action Decoder's public-belief formulation
\cite{foerster2019bad}, its factorised-prescription successor
\cite{sokota2021capi}, and the learned belief models above --- is aimed at games
whose exact belief is intractable, which is the opposite of the situation here.
And the exact-counting machinery this study does not need to approximate ---
Ryser's inclusion--exclusion formula \cite{ryser}, Glynn's sign-vector formula
\cite{glynn}, the Jerrum--Sinclair--Vigoda chain \cite{jsv}, matrix scaling
\cite{lsw} and sequential importance sampling \cite{chen-sis} --- is recorded in
\S\ref{app:rel-counting}, along with why the block structure of this game
sidesteps all of it. History filtering is FNP-complete in general
\cite{solinas-history}; Fish escapes because its public record determines every
card movement, so the only object to reconstruct is the initial deal. No learned
belief model appears anywhere in this study, and that is a stated non-goal
rather than an oversight.

\subsection{Team games and common information}
\label{sec:related-team}

Three teammates share a payoff, hold different information and cannot
communicate. The formal statement of that situation is the common-information
approach: reformulate the decentralised problem from the point of view of a
coordinator who knows the common information and chooses \emph{prescriptions},
maps from each controller's private information to its action
\cite{nayyar2013common}. In Fish the common information is the entire public
history, so the coordinator's information state is exactly the belief object the
engine already maintains --- which is the formal content of the mechanism in
\S\ref{sec:mech-knowledge} that replaces six per-seat deduction states with one
shared object plus per-seat refinement. It is also the reason the SPARTA
constraint above exists: if one seat searches and its teammates do not, they are
no longer executing the prescription the first seat's belief update assumes, and
the reduction silently fails. Solving the coordinator's problem is out of reach
--- a prescription maps every possible nine-card hand to an ask --- so this is a
correctness lens, not an algorithm.

The team-game equilibrium literature is cited to delimit scope. The target
concept for a Fish team is a team-maxmin equilibrium with correlation
\cite{celli-gatti}, connected to extensive-form correlation by Farina and
coauthors \cite{farina-collusion} and made computable through the team belief
DAG and the two-player conversion \cite{zhang-tbdag,carminati-tpi}, at a
representation cost exponential in how much private information distinguishes
teammates inside a public state. Three unrestricted nine-card hands is exactly
the case that cost forbids. Hidden-role machinery \cite{carminati2024hiddenrole}
does not apply at all, because Fish teams are public from the deal. Nothing in
this report is an equilibrium result of any kind.

Conventions are the other half of team play, and the coordination literature
supplies both the vocabulary and the warning. Hanabi is the reference benchmark
for coordination through observable actions \cite{bard-hanabi}; Other-Play
demonstrated how much of a self-play score is convention rather than skill
\cite{hu-otherplay}, the Simplified Action Decoder showed what conventions need
in order to survive exploratory training \cite{hu-sad}, and Off-Belief Learning
supplies an explicit dial on convention depth together with grounded play at
level one \cite{hu2021obl}. Grounded play is unusually strong in this game
because ask legality is itself a hard certificate: an ask proves its asker holds
another card of the named half-suit, so reading a teammate's ask grounded is
exact rather than approximate. Self-explaining deviations formalise the opposite
move --- deviate so that a teammate's theory of mind concludes circumstances are
abnormal --- and produced the first algorithmic finesse plays in Hanabi
\cite{hu2022sed}; in Fish the same channel is supplied by the rules rather than
learned, which is a structural advantage and, in cross-play, a structural risk.
Joint policy search is the bridge-bidding analogue, scoring simultaneous changes
at several information sets, at $+0.63$ IMPs per board against a strong
commercial program \cite{tian2020jps}; the Hanabi construction that most clearly
fails to port is the hat-guessing code, which works because each receiver sees
every other receiver's hand \cite{cox-hanabi}, and in Fish no player sees any
hand.

Human-regularised search is where that risk is priced. piKL anchors a search
result to a reference policy with a bounded divergence \cite{jacob2022pikl}, and
the same construction was extended to coordination with real human partners
\cite{hu2022humanregularized} and used as the planning core of a Diplomacy agent
whose game family --- adversarial teams acting through public actions --- is the
one Fish belongs to \cite{fair2022cicero}. The evidence that a small human
corpus is not the way in is also published: in the ad-hoc human-AI coordination
challenge, an agent trained with no human data at all outperformed a
best-response-to-behavioural-cloning baseline \cite{dizdarevic2025ah2ac2}. This
study uses none of these as method. It takes from them the requirement that the
partner regime be a stated parameter of the policy rather than an assumption,
which is \S\ref{sec:mech-partner}, and that a convention-bearing mechanism be
reported with its cross-play cell and not only its self-play cell.

The last entry is a counterweight to the rest of this section. A 2026
gold-standard study of lightweight game-playing agents concludes that the
binding constraint on strength is information rather than model capacity, and
that disciplined evaluation against a held-out expert beats deeper search
\cite{kelidari2026goldstandard}. That is this study's own thesis stated by
someone else, and it is the reason two of this report's four principal results
are negative and are reported as such.
